\documentclass[10pt,a4paper]{amsart}
\usepackage[margin=19mm]{geometry}
\usepackage{amsmath,amssymb,mathtools}
\usepackage[scr=rsfs,cal=euler]{mathalfa}
\usepackage{xfrac}
\usepackage[unicode,hidelinks,bookmarks=false]{hyperref}
\newcommand{\ie}{i.e.\ }
\newcommand{\cf}{cf.\ }
\newcommand{\resp}{respectively\ }
\newcommand{\Subsection}{\subsection}
\providecommand{\nobreakdash}{\nobreak\hbox{-}}
\setlength{\emergencystretch}{2em}
\multlinegap0pt
\usepackage{amssymb}
\usepackage{enumerate}
\newtheorem{lemma}{Lemma}
\title{Some properties of sets, functions, and multi-objective optimization problems using p-convexity}
\author{Cristian Vera Donoso}
\date{Source: arXiv:2604.09882v1 (CC BY 4.0)}
\begin{document}
\maketitle

\noindent\emph{Source and licence.} Some properties of sets, functions, and multi-objective optimization problems using p-convexity. Version arXiv:2604.09882v1 (CC BY 4.0), distributed by arXiv under the Creative Commons Attribution 4.0 International licence, http://creativecommons.org/licenses/by/4.0/, as recorded in the arXiv licence metadata for this exact version and shown on the abstract page https://arxiv.org/abs/2604.09882v1. Full licence notice: \texttt{licenses/arxiv-2604.09882-CC-BY-4.0.txt}.\par\smallskip

\noindent\textit{Editorial scope.} Counted unit: the article's lemma convexo (source0.tex lines 438--445). The article's complete original proof of it is delivered with it (source0.tex lines 447--507, one \texttt{proof} environment of 61 lines). No mathematics is written, repaired or re-derived: the statement and the proof are the article's own lines, in the article's own notation, and no retained line is substituted or re-flowed.  No further source-local context is needed: every cross-reference of every retained line resolves inside the two delivered spans.  Nothing was omitted from any retained span, so the package carries no omission record.  No retained line contains a citation, so the wrapper carries no bibliography and no bibliography entry.  No retained line contains a figure environment, an image-inclusion command or drawing code, so no image is delivered and the package declares no asset dependency.  Editorial formatting only, in addition: an AMS \texttt{amsart} preamble that restates the article's own declarations and macros used by the retained lines, namely the environments \texttt{lemma} on separate counters, together with the standard packages those lines need.  The retained environments print in this document as Lemma 1 in order of appearance, whereas the article prints the same items with the numbering of its own full document.  The complete native source of the article as distributed in the arXiv e-print for this exact version is bundled unchanged as \texttt{sources/198231/source0.tex}.
\begin{lemma}
Let $K \subseteq \mathbb{R}^{n}$ be a $p$-convex set, with
$x \in \operatorname{int} K$ and $y \in \overline{K}$. Then
\[
[x,y)_{p} \subseteq \operatorname{int} K.
\]
\label{trazo p convexo}
\end{lemma}

\begin{proof}
We distinguish two cases.

\medskip
\noindent\textbf{Case 1: $y \in K$.}
Let $z \in [x,y)_{p}$. Then there exists $\lambda \in (0,1]$ such that
\[
z = \lambda x + (1-\lambda^{p})^{1/p} y.
\]
Since $x \in \operatorname{int} K$, there exists $\delta > 0$ such that
$\mathbb{B}_{q}(x,\delta) \subseteq K$.
We claim that $\mathbb{B}_{q}(z,\lambda\delta) \subseteq K$.
Indeed, let $u \in \mathbb{B}_{q}(z,\lambda\delta)
= z + \lambda \mathbb{B}_{q}(0,\delta)$.
Then $u = z + \lambda w$ for some $w \in \mathbb{B}_{q}(0,\delta)$, and hence
\[
u = \lambda(x+w) + (1-\lambda^{p})^{1/p} y.
\]
Since $x+w \in \mathbb{B}_{q}(x,\delta) \subseteq K$, $y \in K$, and $K$ is
$p$-convex, it follows that $u \in K$.
Thus $z \in \operatorname{int} K$, and therefore
\[
[x,y)_{p} \subseteq \operatorname{int} K \quad \text{for all } y \in K.
\]

\medskip
\noindent\textbf{Case 2: $y \in \overline{K} \setminus K$.}
Let
\[
z = \lambda x + (1-\lambda^{p})^{1/p} y
\quad \text{for some } \lambda \in (0,1).
\]
Since $x \in \operatorname{int} K$, there exists $\delta > 0$ such that
$\mathbb{B}_{q}(x,\delta) \subseteq K$.
Because $y \in \overline{K}$, there exists
\[
k \in K \cap \mathbb{B}_{q}\!\left(
y, \frac{\lambda}{(1-\lambda^{p})^{1/p}}\delta \right).
\]
Hence $k = y + \frac{\lambda}{(1-\lambda^{p})^{1/p}} v$ for some
$v \in \mathbb{B}_{q}(0,\delta)$, and therefore
\[
(1-\lambda^{p})^{1/p} y
= (1-\lambda^{p})^{1/p} k - \lambda v.
\]
Substituting into the expression for $z$, we obtain
\[
z = \lambda(x-v) + (1-\lambda^{p})^{1/p} k.
\]
Since $x-v \in \mathbb{B}_{q}(x,\delta) \subseteq K$ and $k \in K$, the
$p$-convexity of $K$ implies that $z \in K$.
Applying Case~1, we conclude that
\[
[x,z)_{p} \subseteq \operatorname{int} K.
\]
As $z \in [x,y)_{p}$ was arbitrary, it follows that
\[
[x,y)_{p} \subseteq \operatorname{int} K.
\]
This completes the proof.
\end{proof}

\end{document}
